\begin{frame}{官方結果（single-pass）}
  \vspace{-0.35em}
  \begin{columns}[T,onlytextwidth]
    \column{0.52\textwidth}
    \footnotesize
    \begin{tabular}{@{}lrrr@{}}
      \toprule
      設定 & \#Corr & \#Ext@30s & PAR@30s \\
      \midrule
      論文參考 & \PaperCorr & \PaperExtThirty & \PaperParThirty s \\
      \textbf{DS t=0.2} & \textbf{\DSZCorr} & \textbf{\DSZExtThirty} & \textbf{\DSZParThirty s} \\
      DS t=0.7 & \DSvCorr & \DSvExtThirty & \DSvParThirty s \\
      \bottomrule
    \end{tabular}
    \column{0.46\textwidth}
    \centering
    \begin{tikzpicture}
\begin{axis}[
  width=0.88\linewidth,
  height=4.2cm,
  xbar,
  xmin=0,
  xmax=65,
  ytick={1,2,3},
  yticklabels={論文參考,DS t=0.2,DS t=0.7},
  nodes near coords,
  nodes near coords align={horizontal},
  bar width=10pt,
  enlarge y limits=0.25,
  xlabel={\#Ext@30s},
]
\addplot[fill=SeriesDS] coordinates {
  (\PaperExtThirty,1)
  (\DSZExtThirty,2)
  (\DSvExtThirty,3)
};
\end{axis}
\end{tikzpicture}

  \end{columns}
  \vspace{0.3em}
  {\footnotesize\textcolor{SeriesPaper}{主設定 t=0.2 · 不宣稱勝負}\par}
  \note{
    \textbf{約 60 秒}

    主設定是 temp 0.2：\#Corr 703，和論文參考 710 接近；\#Ext 在 30 秒是 48，論文參考是 21；PAR 在 30 秒也明顯較低。0.7 那次 \#Ext@30s 更高但 \#Corr 更低，說明 temperature 不是愈高愈好。這些都是 print\_results.py 的 official 定義，我們不把它講成打贏論文。
  }
\end{frame}
